% ============================================================================
\chapter{Compiling to Lookup Tables}
\label{ch:lut}
% ============================================================================

This chapter describes Phase~3 of the hybrid architecture: compiling
any remaining continuous functions into discrete lookup tables (LUTs).
After this step, the entire inference pipeline operates on Boolean
logic, integer arithmetic, and memory reads---with zero floating-point
operations.


% ----------------------------------------------------------------------------
\section{The LUT Compilation Idea}
\label{sec:lut-idea}
% ----------------------------------------------------------------------------

A lookup table replaces computation with memory: instead of evaluating
$f(x)$ at runtime, we precompute $f(x)$ for every possible input $x$
and store the results. At inference time, we simply read $f(x)$ from
memory.

\begin{equation}
  \text{LUT}[i] = f(i) \quad \text{for } i = 0, 1, \ldots, M{-}1
  \label{eq:lut-def}
\end{equation}

This is only practical when the number of possible inputs $M$ is
small. For a B-spline~\cite{deboor1978splines} scoring function whose input is a popcount
value in $[0, 200]$, $M = 201$---trivially small.

\begin{figure}[htbp]
\centering
\begin{tikzpicture}[
  box/.style={draw, rounded corners, minimum width=2cm,
    minimum height=0.8cm, font=\small\sffamily, align=center},
  arr/.style={-{Stealth[length=4pt]}, thick},
]
  % Training time
  \node[font=\sffamily\bfseries] at (3, 3) {Training time};
  \node[box, fill=bitorange!10] (input1) at (0, 2) {popcount\\value};
  \node[box, fill=bitblue!10] (spline) at (3.5, 2) {B-spline\\$\phi_k(x)$};
  \node[box, fill=bitgreen!10] (score1) at (7, 2) {class\\score};
  \draw[arr] (input1) -- (spline);
  \draw[arr] (spline) -- (score1);

  % Compilation arrow
  \draw[->, very thick, bitpurple, dashed] (3.5, 1.2) -- (3.5, 0.3)
    node[midway, right, font=\scriptsize\sffamily] {compile};

  % Inference time
  \node[font=\sffamily\bfseries] at (3, -0.5) {Inference time};
  \node[box, fill=bitorange!10] (input2) at (0, -1.5) {popcount\\value};
  \node[box, fill=bitpurple!10] (lut) at (3.5, -1.5) {LUT\\read};
  \node[box, fill=bitgreen!10] (score2) at (7, -1.5) {class\\score};
  \draw[arr] (input2) -- (lut);
  \draw[arr] (lut) -- (score2);
\end{tikzpicture}
\caption{LUT compilation replaces a B-spline evaluation (training time)
  with a simple table read (inference time). The function is computed
  once and stored; inference is a single memory access.}
\label{fig:lut-compile}
\end{figure}


% ----------------------------------------------------------------------------
\section{The Compilation Process}
\label{sec:compilation}
% ----------------------------------------------------------------------------

For each scoring function $\phi_k$:

\begin{enumerate}
  \item \textbf{Determine the input range.} The input is a popcount
    of binary gate outputs, so $x \in \{0, 1, \ldots, n_k\}$ where
    $n_k$ is the number of output gates for class $k$.

  \item \textbf{Evaluate the trained B-spline} at every integer
    input: compute $\phi_k(0), \phi_k(1), \ldots, \phi_k(n_k)$ in
    floating point.

  \item \textbf{Quantise to integers.} Convert each floating-point
    score to a fixed-point integer representation:
    \begin{equation}
      \text{LUT}_k[i] = \text{round}\!\left(
        \frac{\phi_k(i) - \phi_{\min}}{\phi_{\max} - \phi_{\min}}
        \times (2^b - 1)\right)
      \label{eq:quantise}
    \end{equation}
    where $b$ is the bit width (e.g., $b = 8$ for 8-bit integers)
    and $\phi_{\min}, \phi_{\max}$ are the range of the function.

  \item \textbf{Store the table.} The result is an array of $n_k + 1$
    integer values.
\end{enumerate}

\begin{engineernote}
With 8-bit quantisation and 201 entries per class:
\[
  \text{Total LUT size} = 10 \times 201 \times 1 \text{ byte}
  = 2{,}010 \text{ bytes} \approx 2 \text{ kB}
\]
This is negligible. Even with 16-bit entries, the total is 4~kB.
\end{engineernote}


% ----------------------------------------------------------------------------
\section{Quantisation Error}
\label{sec:quant-error}
% ----------------------------------------------------------------------------

Quantising a continuous function to $b$-bit integers introduces
\concept{quantisation error}. The maximum error per entry is:

\begin{equation}
  \epsilon_{\max} = \frac{\phi_{\max} - \phi_{\min}}{2^{b+1}}
  \label{eq:quant-error}
\end{equation}

For $b = 8$: $\epsilon_{\max} = \frac{\phi_{\max} - \phi_{\min}}{512}$.
Since the scoring function's range is typically small (e.g., $[-5, 5]$
for log-probability scores), the quantisation error is about 0.02 per
entry. This rarely changes the argmax classification.

\begin{keyinsight}
The classification decision depends on the \emph{relative ordering}
of class scores, not their absolute values. Quantisation preserves
relative ordering as long as the quantisation error is smaller than
the gap between the top two class scores. For well-trained models,
this gap is usually large (the correct class is confident), so 8-bit
quantisation introduces negligible accuracy loss.
\end{keyinsight}


% ----------------------------------------------------------------------------
\section{ULEEN: An Extreme LUT Architecture}
\label{sec:uleen}
% ----------------------------------------------------------------------------

The ULEEN architecture~\cite{susskind2024uleen} pushes the LUT idea
to its logical extreme: \emph{the entire network is lookup tables.}

In ULEEN:
\begin{enumerate}
  \item The input bits are grouped into small subsets (e.g., 8 bits
    each).
  \item Each subset indexes into a LUT with $2^8 = 256$ entries.
  \item The LUT outputs are combined (summed or concatenated) and
    fed into the next layer of LUTs.
  \item This continues for 2--3 layers.
\end{enumerate}

ULEEN achieves 96.2\% accuracy on MNIST~\cite{lecun1998mnist} with a model size of only
16.9~kB. Inference requires only memory reads and integer
additions---no multiplications, no logic gate evaluations, not even
comparisons (beyond the final argmax).

\begin{table}[htbp]
\centering
\caption{Comparison of LUT-based approaches.}
\label{tab:lut-comparison}
\begin{tabular}{@{}lccc@{}}
  \toprule
  \textbf{Approach} & \textbf{Accuracy} & \textbf{Model size}
    & \textbf{Operations} \\
  \midrule
  ULEEN~\cite{susskind2024uleen} & 96.2\% & 16.9 kB & Table reads + adds \\
  Our hybrid (projected) & 97--98\% & $\sim$1 MB & AND + gate LUT + score LUT \\
  2-layer FP network & 98.0\% & $\sim$940 kB & FP multiply-add \\
  \bottomrule
\end{tabular}
\end{table}

\begin{engineernote}
ULEEN shows that pure-LUT architectures are viable. Our hybrid
approach trades a larger model size for higher accuracy by using the
Tsetlin Machine and LGN to build a richer feature representation
before the final LUT scoring.
\end{engineernote}


% ----------------------------------------------------------------------------
\section{The Complete Inference Cost}
\label{sec:inference-cost}
% ----------------------------------------------------------------------------

After LUT compilation, the total inference cost for one MNIST image:

\begin{table}[htbp]
\centering
\caption{Operation count for the complete hybrid inference pipeline.}
\label{tab:inference-cost}
\begin{tabular}{@{}llr@{}}
  \toprule
  \textbf{Stage} & \textbf{Operation type} & \textbf{Count} \\
  \midrule
  Booleanise + expand & Comparisons, NOT & $\sim$1,600 \\
  TM clause evaluation & Bitwise AND, compare & $\sim$250,000 \\
  LGN gate evaluation & 4-entry LUT reads & $\sim$16,000 \\
  Popcount & Integer addition & $\sim$5,000 \\
  Score lookup & 201-entry LUT reads & 10 \\
  Argmax & Integer comparisons & 9 \\
  \midrule
  \textbf{Total} & & $\sim$273,000 \\
  \bottomrule
\end{tabular}
\end{table}

All operations are: bitwise logic, integer comparison, integer
addition, and memory reads from small tables. The energy cost per
operation is 0.1--5~pJ (logic and SRAM), compared to 3.7~pJ for each
floating-point multiply in a conventional network~\cite{horowitz2014energy}.

\begin{keyinsight}
The hybrid architecture achieves its energy efficiency not by having
fewer operations, but by using \emph{cheaper} operations. Each
operation costs 10--40$\times$ less energy than a floating-point
multiply-accumulate, and the model fits entirely in L1/L2 cache,
avoiding the 6,400$\times$ cost of DRAM accesses~\cite{horowitz2014energy}.
\end{keyinsight}


% ----------------------------------------------------------------------------
\section{Generalisation: LUTs for Any Bounded-Input Function}
\label{sec:generalisation}
% ----------------------------------------------------------------------------

The LUT compilation technique applies whenever:
\begin{enumerate}
  \item The input to the function has a small, finite number of
    possible values.
  \item The function is known (trained) before deployment.
  \item There is enough memory to store the table.
\end{enumerate}

In our architecture, the scoring functions have at most 201 possible
inputs (popcount values 0--200). But the same technique could be
applied to:

\begin{itemize}
  \item \textbf{Non-linear combination of two integer inputs:}
    If two popcounts $a \in [0, 50]$ and $b \in [0, 50]$ need to be
    combined, a 2D LUT of $51 \times 51 = 2601$ entries stores the
    result.
  \item \textbf{Activation functions in quantised networks~\cite{hubara2017quantized}:} A
    quantised activation with 8-bit input has $256$ possible values,
    all storable in a 256-entry table.
  \item \textbf{Distance computations:} If both inputs are bounded
    integers, Hamming distance or other metrics can be precomputed.
\end{itemize}

The cost of a LUT scales as $\prod_i |D_i|$ where $|D_i|$ is the
number of possible values for input $i$. For one input with 256
values: 256 entries. For two inputs with 256 values each:
$256^2 = 65{,}536$ entries. For three: $256^3 \approx 16.7$ million.
The exponential growth limits LUTs to functions of few inputs, which
is why we use them only for the final scoring layer.


% ----------------------------------------------------------------------------
\section{Summary}
\label{sec:ch11-summary}
% ----------------------------------------------------------------------------

LUT compilation is the final step that eliminates all floating-point
arithmetic from the inference pipeline:

\begin{enumerate}
  \item Trained B-spline scoring functions are evaluated at every
    possible integer input.
  \item The results are quantised to fixed-precision integers and
    stored in small tables.
  \item At inference time, a table read replaces the function
    evaluation.
  \item The quantisation error is negligible for well-trained models
    with 8-bit precision.
\end{enumerate}

With this step complete, we have all three components of the hybrid
architecture. Part~V assembles them into a single system and analyses
the result.
